% -*- coding: utf-8 -*-
%%%%%%%%%%%%%%%%%%%%%%%%%%%%%%%%%%%%%%%%%%%%%%%%%%%%%%%%%%%%%%%%%%%%%%%%% 
% 改用 ctex 宏包
% Last Update: 2013-12-31
%%%%%%%%%%%%%%%%%%%%%%%%%%%%%%%%%%%%%%%%%%%%%%%%%%%%%%%%%%%%%%%%%%%%%%%%%
%
%   LaTeX File for Master Thesis of Nanjing University
%   LaTeX + CJK     南京大学博士（硕士）论文模板
%   Based on Shangzhi Ye Template for Tsinghua
%   Version: 1.0
%   Last Update: 2005-3-31
%   Author: Guanghong Zuo
%
%%%%%%%%%%%%%%%%%%%%%%%%%%%%%%%%%%%%%%%%%%%%%%%%%%%%%%%%%%%%%%%%%%%%%%%%%
%
%   LaTeX File for Master Thesis of Tsinghua University
%   LaTeX + CJK     清华大学硕士论文模板
%   Based on Wang Lei's Template for Tsinghua
%   Version: 1.0
%   Last Update: 2004-12-25
%   Author: Shaozhi Ye
%
%%%%%%%%%%%%%%%%%%%%%%%%%%%%%%%%%%%%%%%%%%%%%%%%%%%%%%%%%%%%%%%%%%%%%%%%%
%
%   LaTeX File for Doctor (Master) Thesis of Tsinghua University
%   LaTeX + CJK     清华大学博士（硕士）论文模板
%   Based on Wang Tianshu's Template for XJTU
%   Version: 1.75
%   Last Update: 2004-04-12
%
%%%%%%%%%%%%%%%%%%%%%%%%%%%%%%%%%%%%%%%%%%%%%%%%%%%%%%%%%%%%%%%%%%%%%%%%%
%   Copyright 2002-2004  by  Lei Wang (BaconChina)       (bcpub@sina.com)
%%%%%%%%%%%%%%%%%%%%%%%%%%%%%%%%%%%%%%%%%%%%%%%%%%%%%%%%%%%%%%%%%%%%%%%%%
%
%   LaTeX File for phd thesis of xi'an Jiao Tong University
%
%%%%%%%%%%%%%%%%%%%%%%%%%%%%%%%%%%%%%%%%%%%%%%%%%%%%%%%%%%%%%%%%%%%%%%%%%
%   Copyright 2002  by  Wang Tianshu    (tswang@asia.com)
%%%%%%%%%%%%%%%%%%%%%%%%%%%%%%%%%%%%%%%%%%%%%%%%%%%%%%%%%%%%%%%%%%%%%%%%%

%\documentclass[a4paper,fancyhdr,fntef,cap,winfonts,hyperref]{ctexbook}
\documentclass[a4paper,fancyhdr,fntef,cap,hyperref]{ctexbook}
\usepackage{shortvrb, makeidx}
% -*- coding: utf-8 -*-
%%%%%%%%%%%%%%%%%%%%%%%%%%%%%%%%%%%%%%%%%%%%%%%%%%%%%%%%%%%
% 引用的宏包和相应的定义
%%%%%%%%%%%%%%%%%%%%%%%%%%%%%%%%%%%%%%%%%%%%%%%%%%%%%%%%%%%
% 图形支持宏包

\usepackage{graphicx}

\usepackage{subfig}
% 支持彩色
\usepackage{color}
% eps图像
\usepackage{epsfig}

% 首行缩进宏包
\usepackage{indentfirst}

% 版面控制宏包，定义规定的版面尺寸
%\usepackage[%paperwidth=18.4cm, paperheight= 26cm,
%            paper=a4paper,
%            vmargin={3.8 cm,3.8 cm},
%            body={14.6true cm,22true cm},
%            %twosideshift=0 pt,
%            %headheight=1.0true cm
%            ]{geometry}

% 中文支持宏包
%\usepackage{CJK}

% 脚注控制
\usepackage[perpage]{footmisc}

% AMSLaTeX宏包 用来排出更加漂亮的公式

\usepackage{amsmath}
\usepackage{amsfonts}
\usepackage{amssymb}
\usepackage{cases}
\usepackage{amstext}
\usepackage{mathtools}
\usepackage{amsthm}
\usepackage{enumitem}
%\usepackage{biblatex}
% 不同于\mathcal or \mathfrak 之类的英文花体字体
\usepackage{mathrsfs}
% 更多数学符号
%\usepackage{mathabx}

% enumerate环境, 控制行距
\setlist{nosep}
% 控制enumerate编号为(1), (2), ...
\renewcommand\labelenumi{(\theenumi)}

% 定理类环境宏包，其中 amsmath 选项用来兼容 AMS LaTeX 的宏包
%\usepackage[amsmath,thmmarks]{ntheorem}

% 因为图形可浮动到当前页的顶部，所以它可能会出现
% 在它所在文本的前面. 要防止这种情况，可使用 flafter
% 宏包
%\usepackage{flafter}

%浮动图形控制宏包
%允许上一个section的浮动图形出现在下一个section的开始部分
%该宏包提供处理浮动对象的 \FloatBarrier 命令，使所有未处
%理的浮动图形立即被处理
\usepackage[below]{placeins}

%在每一断页clearpage主要对浮动图形有用
%% \usepackage{afterpage}
%% \afterpage{\clearpage}

% 图文混排用宏包
\usepackage{floatflt}

% 图形和表格的控制
\usepackage{rotating}

% tex1cm宏包，控制字体的大小
%\usepackage{type1cm}

% 控制标题的宏包
%\usepackage[sf]{titlesec}

% 控制目录的宏包
%\usepackage{titletoc}

% 处理数学公式中的黑斜体的宏包
\usepackage{bm}

%可将浮动对象放置到文件的最后
%\usepackage{endfloat}

% fancyhdr宏包 页眉和页脚的相关定义
\usepackage{fancyhdr}
%\usepackage{fancyref}

%浮动图形和表格标题样式
\usepackage{caption}
\DeclareCaptionLabelSeparator{blank}{\hspace{1em}}

% 定制表格和图形的多行标题行距
\usepackage{setspace}

% 打印当前页面格式的宏包
\usepackage{layouts}


% 使用Times字体的宏包
%\usepackage{times}

% 表格
\usepackage{array}
\usepackage{booktabs}

%双模下划线
%\usepackage[normalem]{ulem}

%关于章节名的一些显示内容，必须放在hyperref之前
%\makeatletter
%\renewcommand\appendix{\par
%  \setcounter{chapter}{0}%
%  \setcounter{section}{0}%
%  \renewcommand\chaptername{\appendixname}
%  \renewcommand\thechapter{\Alph{chapter}}}
%\makeatother

% 生成有书签的pdf及其开关

\iffalse

\def\a{true}
\ifx\a\useyap
\usepackage[ dvipdf,
             CJKbookmarks=true,
             bookmarksnumbered=true,
             bookmarksopen=true,
             colorlinks=true,
             citecolor=blue,
             linkcolor=black,
             anchorcolor=green,
             urlcolor=blue
              ]{hyperref}
        \else
         \usepackage[dvipdfm,
             CJKbookmarks=true,
             bookmarksnumbered=true,
             bookmarksopen=true,
             colorlinks=true,
             citecolor=blue,
             linkcolor=blue,
             anchorcolor=green,
             urlcolor=blue
             ]{hyperref}
             \AtBeginDvi{\special{pdf:tounicode GBK-EUC-UCS2}} % GBK -> Unicode
\fi

\fi

% 支持引用的宏包
%\usepackage{cite}
% 支持引用缩写的宏包
\usepackage{natbib}
\usepackage{hypernat}


\makeindex

\MakeShortVerb{\|}

\topmargin -0.5 true cm
\oddsidemargin 0 true cm
\evensidemargin 0 true cm
\textheight 23 true cm
\textwidth 16.5 true cm

\setlength{\parskip}{0.75ex plus .2ex minus .5ex}
\renewcommand{\baselinestretch}{1.2}

\pagestyle{fancy}

 %\CTEXsetup[name={第,节},number={\chinese{section}},format+={\bfseries},beforeskip={-10ex plus -.1ex minus -.1ex},afterskip={1ex plus .1ex minus .1ex}]{section}
 %\CTEXsetup[name={附录},number={\chinese{chapter}}]{appendix}
 %\CTEXsetup[name={第,章},number={\arabic{chapter}}]{chapter}

 %\setcounter{secnumdepth}{4}
 %\CTEXsetup[name={(,)},number={\arabic{paragraph}}]{paragraph}

% pdf 书签设置
\hypersetup{CJKbookmarks,%
bookmarksnumbered,%
colorlinks,%
linkcolor=black,%
citecolor=black,%
plainpages=false,%
pdfstartview=FitH}

\begin{document}

% 定义所有的eps文件在 figures 子目录下
\graphicspath{{figures/}}

% 文本格式定义
% -*- coding: utf-8 -*-
%%%%%%%%%%%%%%%%%%%%%%%%%%%%%%%%%%%%%%%%%%%%%%%%%%%%%%%%%%%
% 主文档 格式定义
%%%%%%%%%%%%%%%%%%%%%%%%%%%%%%%%%%%%%%%%%%%%%%%%%%%%%%%%%%%

% 将版芯控制在240mm以内, 正文范围控制在220mm以内
% 将版芯控制在240mm以内, 正文范围控制在220mm以内
\addtolength{\headsep}{-10pt}          %页眉位置
%\addtolength{\footskip}{0}         %页脚位置

%%%%%%%%%%%%%%%%%%%%%%%%%%%%%%%%%%%%%%%%%%%%%%%%%%%%%%%%%%
%  页面设置
%%%%%%%%%%%%%%%%%%%%%%%%%%%%%%%%%%%%%%%%%%%%%%%%%%%%%%%%%
%\setlength{\voffset}{-12pt}
%\setlength{\topmargin}{12pt}
%\addtolength{\textheight}{22pt}
%
%\setlength{\hoffset}{15pt}
%\setlength{\oddsidemargin}{20pt}
%\setlength{\evensidemargin}{20pt}

%%%%%%%%%%%%%%%%%%%%%%%%%%%%%%%%%%%%%%%%%%%%%%%%%%%%%%%%%%%
% 公式的精调
%%%%%%%%%%%%%%%%%%%%%%%%%%%%%%%%%%%%%%%%%%%%%%%%%%%%%%%%%%%
%\setlength{\mathindent}{4.7 em}     %左对齐公式缩进量

% \eqnarray如果很长，影响分栏、换行和分页（整块挪动，造成页面空白），
% 可以设置成为自动调整模式
\allowdisplaybreaks[4]

%%%%%%%%%%%%%%%%%%%%%%%%%%%%%%%%%%%%%%%%%%%%%%%%%%%%%%%%%%%
%下面这组命令使浮动对象的缺省值稍微宽松一点，从而防止幅度
%对象占据过多的文本页面，也可以防止在很大空白的浮动页上放置
%很小的图形。
%%%%%%%%%%%%%%%%%%%%%%%%%%%%%%%%%%%%%%%%%%%%%%%%%%%%%%%%%%%
\renewcommand{\textfraction}{0.15}
\renewcommand{\topfraction}{0.85}
\renewcommand{\bottomfraction}{0.65}
\renewcommand{\floatpagefraction}{0.60}


%%%%%%%%%%%%%%%%%%%%%%%%%%%%%%%%%%%%%%%%%%%%%%%%%%%%%%%%%%%
%下面这组命令可以使公式编号随着每开始新的一节而重新开始。
%%%%%%%%%%%%%%%%%%%%%%%%%%%%%%%%%%%%%%%%%%%%%%%%%%%%%%%%%%%
%\makeatletter      % '@' is now a normail "letter" for TeX
%\@addtoreset{eqation}{section}
%\makeatother       % '@' is restored as a "non-letter" character for TeX

%%%%%%%%%%%%%%%%%%%%%%%%%%%%%%%%%%%%%%%%%%%%%%%%%%%%%%%%%%%
% 重定义字体命令
%%%%%%%%%%%%%%%%%%%%%%%%%%%%%%%%%%%%%%%%%%%%%%%%%%%%%%%%%%%
% 注意win2000没有 simsun,　最好到网上找一个
% 一些字体是office2000 带的
%%%%%%%%%%%%%%%%%%%%%%%%%%%%%%%%%%%%%%%%%%%%%%%%%%%%%%%%%%%
%
\newcommand{\song}{\CJKfamily{song}}    % 宋体   (Windows自带simsun.ttf)
\newcommand{\fs}{\CJKfamily{fs}}        % 仿宋体 (Windows自带simfs.ttf)
\newcommand{\kai}{\CJKfamily{kai}}      % 楷体   (Windows自带simkai.ttf)
\newcommand{\hei}{\CJKfamily{hei}}      % 黑体   (Windows自带simhei.ttf)
\newcommand{\li}{\CJKfamily{li}}        % 隶书   (Windows自带simli.ttf)

%%%%%%%%%%%%%%%%%%%%%%%%%%%%%%%%%%%%%%%%%%%%%%%%%%%%%%%%%%%
% 重定义字号命令
%%%%%%%%%%%%%%%%%%%%%%%%%%%%%%%%%%%%%%%%%%%%%%%%%%%%%%%%%%%
%
\newcommand{\yihao}{\zihao{1}}    % 一号, 1.4倍行距
\newcommand{\erhao}{\zihao{2}}    % 二号, 1.25倍行距
\newcommand{\xiaoer}{\zihao{-2}}    % 小二, 单倍行距
\newcommand{\sanhao}{\zihao{3}}    % 三号, 1.5倍行距
\newcommand{\xiaosan}{\zihao{-3}}    % 小三, 1.5倍行距
\newcommand{\sihao}{\zihao{4}}    % 四号, 1.5倍行距
\newcommand{\banxiaosi}{\fontsize{13pt}{19.5pt}\selectfont}    % 半小四, 1.5倍行距
\newcommand{\xiaosi}{\zihao{-4}}    % 小四, 1.5倍行距
\newcommand{\dawuhao}{\fontsize{11pt}{11pt}\selectfont}    % 大五号, 单倍行距
\newcommand{\wuhao}{\zihao{5}}    % 五号, 单倍行距


%%%%%%%%%%%%%%%%%%%%%%%%%%%%%%%%%%%%%%%%%%%%%%%%%%%%%%%%%%%
% 重定义一些正文相关标题
%%%%%%%%%%%%%%%%%%%%%%%%%%%%%%%%%%%%%%%%%%%%%%%%%%%%%%%%%%%

%\theoremstyle{plain} \theorembodyfont{\song\rmfamily}
%\theoremheaderfont{\hei\rmfamily} \theoremseparator{:}
%
%\renewcommand\contentsname{目~~~~~~录}
%\renewcommand\listfigurename{插图}
%\renewcommand\listtablename{表格}
%\renewcommand\indexname{索引}
%\renewcommand\figurename{图}
%\renewcommand\tablename{表}

%\newtheorem{definition}{\hei 定义}[chapter]
%\newtheorem{proposition}[definition]{\hei 命题}
%\newtheorem{lemma}[definition]{\hei 引理}
%\newtheorem{theorem}{\hei 定理}[chapter]
%\newtheorem{axiom}{\hei 公理}
%\newtheorem{corollary}[definition]{\hei 推论}
%\newtheorem{exercise}[definition]{}

%\theoremheaderfont{\CJKfamily{hei}\rmfamily}\theorembodyfont{\rmfamily}
%\theoremstyle{nonumberplain} \theoremseparator{:}
%\theoremsymbol{$\blacksquare$}
%\newtheorem{proof}{\hei 证明}

%\theoremsymbol{$\square$}
%\newtheorem{example}{\hei 例}


%%%%%%%%%%%%%%%%%%%%%%%%%%%%%%%%%%%%%%%%%%%%%%%%%%%%%%%
% 定义页眉和页脚 使用fancyhdr 宏包
%%%%%%%%%%%%%%%%%%%%%%%%%%%%%%%%%%%%%%%%%%%%%%%%%%%%%%%%
%\newcommand{\makeheadrule}{%
%    \makebox[0pt][l]{\rule[.7\baselineskip]{\headwidth}{0.8pt}}%
%    \vskip-.8\baselineskip}%
%
%  \makeatletter
%  \renewcommand{\headrule}{%
%    {\if@fancyplain\let\headrulewidth\plainheadrulewidth\fi
%      \makeheadrule}}%
%
%  \pagestyle{fancy}
%
%去掉章节标题中的数字
%\renewcommand{\chaptermark}[1]{\markboth{\chaptername \ #1}{}}

%\fancyhf{}
%\fancyhead[CO]{\wuhao\leftmark}
%\fancyhead[CE]{\wuhao\leftmark}
%\fancyfoot[C,C]{\wuhao\thepage}

%%%%%%%%%%%%%%%%%%%%%%%%%%%%%%%%%%%%%%%%%%%%%%%%%%%%%%%%
% 设置行距和段落间垂直距离
%%%%%%%%%%%%%%%%%%%%%%%%%%%%%%%%%%%%%%%%%%%%%%%%%%%%%%%%

% 段落之间的竖直距离
\setlength{\parskip}{3pt plus1pt minus1pt}

% 定义行距
\renewcommand{\baselinestretch}{1.2}


%%%%%%%%%%%%%%%%%%%%%%%%%%%%%%%%%%%%%%%%%%%%%%%%%%%%%%%
% 修改引用的格式，
%%%%%%%%%%%%%%%%%%%%%%%%%%%%%%%%%%%%%%%%%%%%%%%%%%%%%%%
% 增加 \ucite 命令使显示的引用为上标形式
\newcommand{\ucite}[1]{$^{\mbox{\scriptsize \cite{#1}}}$}

%%%%%%%%%%%%%%%%%%%%%%%%%%%%%%%%%%%%%%%%%%%%%%%%%%%%%%%%%%%
% 定制浮动图形和表格标题样式
%%%%%%%%%%%%%%%%%%%%%%%%%%%%%%%%%%%%%%%%%%%%%%%%%%%%%%%%%%%
%添加新的表格控制符，用法与p{xx}相似
%P{}表示定制表格例的宽度并且使不够长度的居中与同等其他列上对齐
%M{}表示定制表格例的宽度并且使不够长度的居中与同等其他列中对齐
%B{}表示定制表格例的宽度并且使不够长度的居中与同等其他列下对齐
%\newcommand{\PreserveBackslash}[1]{\let\temp=\\#1\let\\=\temp}
%\newcolumntype{P}[1]{>{\PreserveBackslash\centering}p{#1}}
%\newcolumntype{M}[1]{>{\PreserveBackslash\centering}m{#1}}
%\newcolumntype{B}[1]{>{\PreserveBackslash\centering}b{#1}}

\renewcommand{\captionfont}{\CJKfamily{song}\rmfamily}
\renewcommand{\captionlabelfont}{\CJKfamily{song}\rmfamily}

% 按清华标准, 去掉图表号后面的:
%不太清楚南大的标准，所以没改
\captionsetup{labelsep=blank}

% 按清华标准, 图表标题字体为11pt, 这里写作大五号
\renewcommand{\captionfont}{\dawuhao}

% 设置Caption的宽度
\setcaptionmargin{.5in}

% define graph scale
\def\gs{1.0}

%%%%%%%%%%%%%%%%%%%%%%%%%%%%%%%%%%%%%%%%%%%%%%%%%%%%%%%%%%%%%%%%%%%%%%
% 自定义项目列表标签及格式 \begin{denselist} 列表项 \end{denselist}
%%%%%%%%%%%%%%%%%%%%%%%%%%%%%%%%%%%%%%%%%%%%%%%%%%%%%%%%%%%%%%%%%%%%%%
\newcounter{newlist} %自定义新计数器
\newenvironment{denselist}[1][可改变的列表题目]{%%%%%定义新环境
\begin{list}{\textbf{\hei #1} \arabic{newlist}：} %%标签格式
    {
    \usecounter{newlist}
     \setlength{\labelwidth}{22pt} %标签盒子宽度
     \setlength{\labelsep}{0cm} %标签与列表文本距离
     \setlength{\leftmargin}{0cm} %左右边界
     \setlength{\rightmargin}{0cm}
     \setlength{\parsep}{0ex} %段落间距
     \setlength{\itemsep}{0ex} %标签间距
     \setlength{\itemindent}{44pt} %标签缩进量
     \setlength{\listparindent}{22pt} %段落缩进量
    }}
{\end{list}}%%%%%

\def\defaultfont{\renewcommand{\baselinestretch}{1.5}
\fontsize{12pt}{13pt}\selectfont}

%%%%%%%%%%%%%%%%%%%%%%%%%%%%%%%%%%%%%%%%%%%%%%%%%%%%%%%%%%%%%%%%%%%%%%
% 封面、摘要、与目录格式定义 --- by Guanghong Zuo Feng Hua & Lei Wang
%%%%%%%%%%%%%%%%%%%%%%%%%%%%%%%%%%%%%%%%%%%%%%%%%%%%%%%%%%%%%%%%%%%%%%
\makeatletter
\def\ctitle#1{\def\@ctitle{#1}}\def\@ctitle{}
\def\cdegree#1{\def\@cdegree{#1}}\def\@cdegree{}
\def\cdepart#1{\def\@cdepart{#1}}\def\@cdepart{}
\def\csubject#1{\def\@csubject{#1}}\def\@csubject{}
\def\cdirection#1{\def\@cdirection{#1}}\def\@cdirection{}
\def\cauthor#1{\def\@cauthor{#1}}\def\@cauthor{}
\def\csupervisor#1{\def\@csupervisor{#1}}\def\@csupervisor{}
\def\stitle#1{\def\@stitle{#1}}\def\@stitle{}
\def\cdate#1{\def\@cdate{#1}}\def\@cdate{}
\long\def\cabstract#1{\long\def\@cabstract{#1}}\long\def\@cabstract{}
\def\ckeywords#1{\def\@ckeywords{#1}}\def\@ckeywords{}

\def\etitle#1{\def\@etitle{#1}}\def\@etitle{}
\def\edegree#1{\def\@edegree{#1}}\def\@edegree{}
\def\edepart#1{\def\@edepart{#1}}\def\@edepart{}
\def\esubject#1{\def\@esubject{#1}}\def\@esubject{}
\def\edirection#1{\def\@edirection{#1}}\def\@edirection{}
\def\eauthor#1{\def\@eauthor{#1}}\def\@eauthor{}
\def\esupervisor#1{\def\@esupervisor{#1}}\def\@esupervisor{}
\def\estitle#1{\def\@estitle{#1}}\def\@estitle{}
\def\eaddress#1{\def\@eaddress{#1}}\def\@eaddress{}
\def\edate#1{\def\@edate{#1}}\def\@edate{}
\long\def\eabstract#1{\long\def\@eabstract{#1}}\long\def\@eabstract{}
\def\ekeywords#1{\def\@ekeywords{#1}}\def\@ekeywords{}

\def\ngrade#1{\def\@ngrade{#1}}\def\@ngrade{}
\def\snumber#1{\def\@snumber{#1}}\def\@snumber{}

\sloppy
\pagenumbering{roman}

\def\makecover{
  \begin{titlepage}
\includegraphics[width=2.5cm]{njuxh}\\
    \parbox[t][.5cm][t]{\textwidth}{
    \begin{center}  \end{center} }
  \begin{center}

      \parbox[t][5.2cm][t]{\textwidth}{\begin{center} \song
          \includegraphics[width=5.0cm]{njuwz}\\
         {\bf\yihao 本~~~科~~~毕~~~业~~~论~~~文\\}
         \end{center}}

    \parbox[t][3.5cm][t]{\textwidth}{
    \begin{center}  \end{center} }


    \parbox[t][6cm][c]{\textwidth}{\begin{center}
      \fontsize{16pt}{8pt}\selectfont \song
      \begin{tabular}{p{4.5em} p{2em} p{4em}@{} p{2.5em} p{6em}}
    院~~~~~~~~系 & \multicolumn{4}{c}{\@cdepart}\\
    \cmidrule(r){2-5}\addlinespace[10pt]
    专~~~~~~~~业 & \multicolumn{4}{c}{\@csubject}\\
    \cmidrule(r){2-5}\addlinespace[10pt]
    题~~~~~~~~目 &  \multicolumn{4}{c}{\@ctitle}\\
    \cmidrule(r){2-5}\addlinespace[10pt]
    年~~~~~~~~级 &  \multicolumn{2}{c}{\@ngrade} & 学号 & \@snumber\\
    \cmidrule(r){2-3} \cmidrule(r){5-5}\addlinespace[10pt]
    学生姓名 & \multicolumn{4}{c}{\@cauthor}\\
    \cmidrule(r){2-5}\addlinespace[10pt]
    指导教师 & \multicolumn{2}{c}{\@csupervisor} & 职称 & \@stitle\\
    \cmidrule(r){2-3} \cmidrule(r){5-5}\addlinespace[10pt]
    \multicolumn{2}{l}{论文提交日期} & \multicolumn{3}{c}{\@cdate}\\
    \cmidrule(r){3-5}
  \end{tabular}\end{center}}

    \parbox[t][2.5cm][t]{\textwidth}{
    \begin{center}  \end{center} }

    %\parbox[t][1cm][b]{\textwidth}{
    %\begin{center} {\sanhao \kai \@cdate} \end{center} }

  \end{center}

    % 封二
    \newpage

    \thispagestyle{empty}

    \parbox[t][16cm][t]{\textwidth}{
    \begin{center}  \end{center} }

    \parbox[b][10em][b]{\textwidth}{\kai \sanhao
      \begin{tabular}{l p{12em}}
        学~~~~~~~~~~~~~~~~号：& {\bf ~~~~\@snumber}\\
        论文答辩日期：& ~~~~~~~~~~年~~~~~~~月~~~~~~~日\\
        指~~~导~~~教~~~师： &~~~~~~~~~~~~~~~~~~~~~~（签字）\\
        \end{tabular}
    }

    % 英文封面
    \newpage

    \thispagestyle{empty}
    \begin{center}
    \parbox[t][1cm][t]{\textwidth}{
    \begin{center}  \end{center} }

    \parbox[t][1cm][c]{\textwidth}{\sanhao
    \begin{center} {\textbf{\textsf Dissertation for {\@edegree}}}\end{center} }

    \parbox[t][4cm][c]{\textwidth}{\erhao
    \begin{center} {\textbf{\textsf{\@etitle}}}\end{center} }

    \parbox[t][4cm][t]{\textwidth}{
    \begin{center}  \end{center} }

    \parbox[t][5cm][b]{\textwidth}{
    \begin{center} {\sanhao \textsf{\textbf{\@eauthor\\
            Supervised by ~\@esupervisor \\ \@esubject\\} }} \end{center} }

    \parbox[t][3cm][c]{\textwidth}{ {\sanhao
    \begin{center}
      \@eaddress
    \end{center} } }

    \parbox[t][2cm][b]{\textwidth}{
    \begin{center} {\sanhao \textbf{\textsf{\@edate}}} \end{center} }
    \end{center}
    \end{titlepage}

    \clearpage

    %Abstract%%%%%%%%%%%%%%%%%%%%%%%%%%%%%%%%%%%%%%%%%%%%%%%%%%%%%
    \defaultfont
    \pagestyle{plain}

      \phantomsection
      \begin{center} \kai \xiaoer
        \uuline{南京大学本科生毕业论文(设计)中文摘要}

        \vspace{15pt}

        {\fontsize{16pt}{8pt}\selectfont
          \begin{tabular}{p{3.5em}@{} p{2.5em} p{2em}@{} p{2em}@{} p{2.5em} p{2em}@{} p{7em}@{} p{3.5em}}
            \multicolumn{3}{l}{毕业论文题目:} &\multicolumn{5}{b{17em}}{\@ctitle}\\
            \cmidrule(r){4-8}\addlinespace[8pt]
            \@cdepart &院系& \multicolumn{2}{c}{\@csubject} &专业& \@ngrade &级本科生姓名:& \@cauthor\\
            \cmidrule(r){1-1}\cmidrule(r){3-4}\cmidrule(r){6-6}\cmidrule(r){8-8}\addlinespace[8pt]
            \multicolumn{4}{c}{指导教师(姓名、职称):} &\multicolumn{2}{c}{\@csupervisor} &\@stitle\\
            \cmidrule{5-8}
          \end{tabular}}

        \vspace{1.5cm}

      % {\hei \sanhao 中文摘要}
      \end{center}
      \addcontentsline{toc}{chapter}{中文摘要}
      \markboth{}{}
      \setcounter{page}{1}
      {\kai \sanhao 摘要：} \\

      \@cabstract\\
      \\
     {\kai \sanhao 关键词：} ~\@ckeywords
 
      \defaultfont
      \clearpage

      \newpage
      \phantomsection
      \begin{center} \kai \xiaoer
        \uuline{南京大学本科生毕业论文(设计)英文摘要}
      \end{center}

      \vspace{-10pt}

      {\xiaosan \parindent 0pt

        \begin{tabular}{@{}l p{22em}}
          THESIS: & \fontsize{16pt}{10pt}\selectfont \@etitle\\
        \end{tabular}

        DEPARTMENT:~~~~\@edepart
        \vspace{-5pt}

        SPECIALIZATION:~~~~\@esubject
        \vspace{5pt}

        UNDERGRADUATE:~~~~\@eauthor
        \vspace{5pt}

        MENTOR:~~~~\@esupervisor
        \vspace{5pt}

        ABSTRACT:~~~~
      }

      \vspace{1cm}

      %\begin{center}{\bf \sanhao Abstract}\end{center}
      \addcontentsline{toc}{chapter}{英文摘要}
      \markboth{}{}

      \@eabstract\\
      \\
      KEYWORDS:~\@ekeywords

      \clearpage

    % Contents %%%%%%%%%%%%%%%%%%%%%%%%%%%%%%%%%%%%%%%%%%%%%%%%%%%%%%
    %\renewcommand{\baselinestretch}{1.5}
    %\fontsize{12pt}{12pt}\selectfont
    %\pagestyle{fancy}

    % \phantomsection
    % \addcontentsline{toc}{chapter}{\contentsname}
    \tableofcontents
    \markboth{\contentsname}{\contentsname}
    \clearpage
    %\defaultfont
    \pagenumbering{arabic}
    \setcounter{page}{1}

}

\makeatother


% 定理环境定义
% -*- coding: utf-8 -*-
\newtheorem{theorem}{Theorem}[section]
\newtheorem{thm}[theorem]{定理} % 中文定理
\newtheorem{lemma}[theorem]{引理}
\newtheorem{definition}[theorem]{定义}
\newtheorem{proposition}[theorem]{命题}
\newtheorem{corollary}[theorem]{推论}
\newtheorem{condition}[theorem]{条件}


% 正文部分
%--- Preface ------------------------

%封面 (含中文摘要, 英文摘要与目录)
% -*- coding: utf-8 -*-
%%%%%%%%%%%%%%%%%%%%%%%%%%%%%%%%%%%%%%%%%%%%%%%%%%%%%%%%%%%%%%%%%%%%%%%%%
% 封面、摘要的相关内容
%%%%%%%%%%%%%%%%%%%%%%%%%%%%%%%%%%%%%%%%%%%%%%%%%%%%%%%%%%%%%%%%%%%%%%%%%

%%%------------中文个人信息------------------------
%\cdegree{学士}                          %学士/硕士/博士
\cdepart{数学系}                      %院系
\csubject{统计学}                   %专业
\ctitle{仿紧空间的性质与刻画}       %题目
%\cdirection{数学分析}                 %研究方向
\cauthor{单铭铭}                        %姓名
\csupervisor{师维学}          %导师姓名
\stitle{教授}               %职称
\cdate{二零一七年五月}                  %中文日期

%--------------英文个人信息------------------------

\etitle{Properties and Descriptions of Paracompact Spaces}%标题
\edegree{Bachelor of Science}          % 学位
\edepart{Department of Mathematics}    %院系
\esubject{Statistics}                % 专业
\edirection{General Topology}         % 研究方向
\eauthor{Mit Shan}              % 姓名
\esupervisor{Professor Weixue Shi}      % 导师姓名
\estitle{Professor}
\eaddress{Department of Mathematics, Nanjing University}     % 英文单位
\edate{May, 2017}                              % 英文日期

\ngrade{2013}                                  % 年级
\snumber{131110066}                           % 学号

%-----------------中文摘要----------------------------
\cabstract{
  本文主要是我对于仿紧拓扑空间的刻画和性质的一个学习总结,
  所有结论均来自教科书和参考资料.
  主要证明在满足$T_3$公理的前提下,几个与拓扑空间仿紧性等价的条件,
  比如其中之一是:拓扑空间$X$是仿紧的,当且仅当$X$的每个开覆盖有$\sigma-$局部有限开加细.
  在完成仿紧性质的基本刻画后,给出几个仿紧空间的性质,使仿紧性不再那么神秘.
}

\ckeywords{仿紧}



%----------------英文摘要-----------------------------
\eabstract{
  This article is my learning summary of some descriptions and properties of paracompact
  spaces. All the conclusions come from textbooks and references.
  Several equivalent conditions of paracompactness under $T_3$ axiom
  will be proved. One of them is: Topological space $X$ is paracompact,
  if and only if every open cover of $X$ has $\sigma$-local-finite open refinement.
  After describing paracompactness, we will show some properties of paracompact
  space, which will uncover the mystery of paracompactness.
}

\ekeywords{paracompact}

\makecover


%符号对照表
%\include{appendix/notation}

%--- body --------------------------

%正文章节
% -*- coding: utf-8 -*-
\chapter{预备及定义}
\label{chap01}

为了方便, 下面我们集中给出本文用的一些定义以及一些记号.

\begin{definition} \label{def:regular space}
  设X为拓扑空间.
  \begin{enumerate}
  \item 若X满足$T_2$公理,则称X为\textbf{Hausdorff}空间.
  \item 若X满足$T_3$公理,则称X为\textbf{正则}空间.
  \item 若X满足$T_4$公理,则称X为\textbf{正规}空间.
  \end{enumerate}
\end{definition}

\begin{definition} \label{def:lindelof}
  拓扑空间X称为\textbf{Lindel\"{o}f}的,如果X的任意开覆盖都有可数子覆盖.
\end{definition}

\begin{definition} \label{def:countably compact}
  拓扑空间X称为\textbf{可数紧}的,如果X的任意可数开覆盖存在有限子覆盖.
\end{definition}

\begin{definition} \label{def:refinement}
  X为一拓扑空间.$\mathscr{U}, \mathscr{V}$为X的两个覆盖.
  称$\mathscr{V}$为$\mathscr{U}$的一个\textbf{加细},如果
  \begin{equation*}
    \forall V \in \mathscr{V}, \exists U \in \mathscr{U}, \text{使得}
    V \subseteq U, \text{且} \newline  %% 这里最好换个行
    \bigcup \mathscr{V} = \bigcup \mathscr{U}.
  \end{equation*}
\end{definition}

\begin{definition} \label{def:local finite}
  X为一拓扑空间,$\mathscr{A}$为X的一个子集族.称$\mathscr{A}$在X中是\textbf{局部有限}的,如果
  \begin{equation*}
    \forall x \in X, \exists x \text{的邻域}U,\text{使得}
    \{A \in \mathscr{A} \mid A \cap U \neq \varnothing \}\text{为有限集}.
  \end{equation*}
\end{definition}

\begin{definition} \label{def:sigma local finit}
  X为一拓扑空间,$\mathscr{A}$为X的一个子集族.
  称$\mathscr{A}$在X中是\textbf{$\sigma$-局部有限}的,如果$\mathscr{A}$是可数多个局部有限族的并.
\end{definition}

\begin{definition} \label{def:discrete}
  拓扑空间X的子集族$\mathscr{U}$称为\textbf{离散的}, 如果对任意$x \in X$,
  存在$x$的邻域$A$, 使得$A$最多与$\mathscr{U}$中的一个元素相交非空.
\end{definition}

\begin{definition} \label{def:closure-preserving}
  设$\mathscr{U}$为拓扑空间X的子集族, 称$\mathscr{U}$为\textbf{闭包保持}的, 如果
  \[
    \bigcup \{ \overline{U} \mid U \in \mathscr{U} \}
    = \overline{\bigcup \{ U \mid U \in \mathscr{U} \}}.
  \]
\end{definition}

%% 仿紧的定义
\begin{definition} \label{def:paracompact}
  拓扑空间X是\textbf{仿紧}的,如果它的任意开覆盖都有局部有限开加细.
\end{definition}
紧空间是仿紧空间,因为紧空间的任意开覆盖存在有限子覆盖,从而是局部有限的,子覆盖也是一个开加细,所以是仿紧的.

\begin{definition}
  设X, Y为拓扑空间,称映射$f : X \longrightarrow Y$为\textbf{完备的},
  如果f为闭的连续满映射,且对$\forall y \in Y, f^{-1}(y)$是X的紧子集.
\end{definition}

% -*- coding: utf-8 -*-
\chapter{仿紧空间的一些性质(上)}
\label{chap02}

在我们得到仿紧空间的定义后, 我们立即就可以得到一些基本的性质.
这些性质并不需要对仿紧空间的深度刻画就可以得到证明.
所以在刻画仿紧空间前, 我们把这些性质列出来, 算是初识仿紧性.

下面的命题说明了仿紧性与紧性之间的相似之处.

\begin{proposition}
  设X是仿紧空间, A是X中的闭集, 则A作为子空间也是仿紧的.
\end{proposition}
\begin{proof}
  设$\mathscr{A}$是$A$在$X$中的一个开覆盖, 则$\mathscr{A} \cup \{ X \setminus A \}$
  是$X$的一个开覆盖, 由$X$的仿紧性, 存在局部有限开加细$\mathscr{L}$.
  令$\mathscr{M} = \{ U \in \mathscr{L} \mid U \cap A \neq \varnothing \}$,
  再令$\mathscr{N} = \{ U \cap A \mid U \in \mathscr{M} \}$,
  则$\mathscr{N}$为$A$作为子空间, 开覆盖$\{ U \cap A \mid U \in \mathscr{A} \}$
  的局部有限开加细. 所以A作为子空间也是仿紧的.
\end{proof}

同样是和紧致性进行类比, 紧致的Hausdorff空间是$T_3$, 我们有下面的定理.

\begin{thm} \label{thm:paracompact T2}
  仿紧Hausdorff空间是$T_3$的.
\end{thm}
\begin{proof}
  设$X$为仿紧Hausdorff空间, $x_0 \in X, F$为$X$中的闭集,且$x_0 \notin F$.
  对$\forall y \in F, \exists U_y \text{为开集且} x_0 \notin \overline{U_y}.$
  则$\{ U_y \mid y \in F\} \cup \{ X \setminus F \}$是$X$的一个开覆盖,由$X$的仿紧性,
  存在局部有限开加细$\mathscr{V} = \{ V_j \mid j \in J \}$.
  存在$x$的开邻域$B$,存在有限集$J^{'} \subseteq J$,使得$\forall j \in J \setminus J^{'},
  V_j \cap B = \varnothing$.对$\forall j \in J^{'}, \exists y_j, \text{使得}
  V_j \subseteq U_{y_j}$.而$x_0 \notin \overline{U_{y_j}}, \text{所以} x_0 \notin \overline{V_j}$.
  我们可以观察到,$\mathscr{V}\text{中的元素要么包含在某个}U_y\text{中},\text{要么包含在} X \setminus F$中.
  令$\mathscr{A} = \{ V_j \in \mathscr{V} \mid V_j \cap F \neq \varnothing \}$,是F的一个开覆盖.
  再由上面的观察,有$x_0 \notin \bigcup \mathscr{A}$.而$\mathscr{A}$中只有有限个元素与B相交,
  且$x_0$不属于这有限个元素的闭包.所以存在$x_0$的开邻域$C \subseteq B$,使得$C \cap \bigcup \mathscr{A}
  = \varnothing$.所以$X$是$T_3$的.
\end{proof}

\begin{lemma} \label{lemma:local finite close union}
  若$\{ U_i \mid i \in I\}$是一个局部有限的子集族,则
  \[
    \bigcup_{i \in I}\overline{U_i} = \overline{\bigcup_{i \in I}U_i} \,\,.
  \]
\end{lemma}
\begin{proof}
  (1) $x \in \bigcup_{i \in I}\overline{U_i} \Longrightarrow x \in \overline{\bigcup_{i \in I}U_i}\,\,$:
  由条件, $\exists i_0 \in I$, 使得$x \in \overline{U_{i_0}}$,
  从而自然的, $x \in \overline{\bigcup_{i \in I}U_i}$.\newline
  (2) $x \in \overline{\bigcup_{i \in I}U_i} \Longrightarrow x \in \bigcup_{i \in I}\overline{U_i}\,\,$:
  (反证)假设$x \notin \bigcup_{i \in I}\overline{U_i}$.
  由局部有限,存在$x$的开邻域$A$,使得对$i \notin \{i_1, i_2, ..., i_N \} \subseteq I$,
  有$A \cap U_i = \varnothing$. 记$I^{'} = \{i_1, i_2, ..., i_N \}$,
  则对$\forall i \in I^{'}, x_0 \notin \overline{U_i}$. 
  令$B = A \cap \overline{U_{i_1}}^c \cap \overline{U_{i_2}}^c \cap \cdots \cap \overline{U_{i_N}}^c$,
  则$B$为非空开集, $x_0 \in B$, 且对$\forall i \in I, B \cap U_i = \varnothing$.
  所以$B \cap \bigcup_{i \in I} U_i = \varnothing$.
  再由$B$是开集, 有$B \cap \overline{\bigcup_{i \in I} U_i} = \varnothing$.
  推出$x \notin \overline{\bigcup_{i \in I} U_i}$, 矛盾.
\end{proof}

这个引理其实就是说, 局部有限集族是闭包保持的. 下面的命题在紧致空间的
对应版本是: 紧致Hausdorff空间是$T_4$的.

\begin{thm}
  仿紧Hausdorff空间是$T_4$的.
\end{thm}
\begin{proof}
  设$X$是仿紧Hausdorff空间,$A$和$B$是$X$中的闭集,且$A \cap B = \varnothing$.
  对$\forall a \in A$,由定理\ref{thm:paracompact T2},
  $\exists$开集$U_a$,使得$\overline{U_a} \cap B = \varnothing$.
  令$\mathscr{A} = \{ U_a \mid a \in A \} \cup \{ X \setminus A \}$,
  则$\mathscr{A}$为$X$的一个开覆盖,由$X$的仿紧性,
  存在局部有限开加细$\mathscr{V} = \{ V_j \mid j \in J \}$.
  令$J^{'} = \{ j \in J \mid V_j \cap A \neq \varnothing \}$,
  则$\mathscr{V}^{'} = \{ V_j \mid j \in J^{'} \}$是$A$的一个开覆盖.
  且对$\forall j \in J^{'}, \exists a \in A$,使得$V_j \subseteq U_a$.
  所以$\overline{V_j} \cap B = \varnothing, \forall j \in J^{'}$.
  因为$\mathscr{V}^{'}$是局部有限的,由引理\ref{lemma:local finite close union},
  $\overline{\bigcup_{j \in J^{'}}V_j} = \bigcup_{j \in J^{'}} \overline{V_j}$.
  所以$B \cap \overline{\bigcup_{j \in J^{'}}V_j} = \varnothing$.
  所以$X$是$T_4$的.
\end{proof}

从上面的性质可以看出, 仿紧性与紧致性是很类似的.
事实上, 仿紧性就是对紧致性放松了一些条件. 仿紧空间不一定是紧的,
比如, 自然数集作为一个离散的拓扑空间, 是仿紧的(通过下面的刻画可以知道).
但是却不是紧的.

% -*- coding: utf-8 -*-
\chapter{$T_3$仿紧空间的一些刻画}
\label{chap03}

接下来, 我们进入本文的正题, 就是对仿紧空间(附加上一些分离公理)的刻画,
使仿紧性质更加清晰一些, 而不仅仅只是个定义.

\begin{thm} \label{thm:paracompact sigma finite}
  设X是$T_3$空间,则下面各条件等价:
  \begin{enumerate}
  \item X是仿紧的;
  \item X的任意开覆盖有$\sigma$-局部有限开加细;
  \item X的任意开覆盖有局部有限加细;
  \item X的任意开覆盖有局部有限闭加细.
  \end{enumerate}
\end{thm}
\begin{proof}
  %% PROOF HERE
  (1) $\Rightarrow$ (2) :
  由仿紧的定义, 是显然的.
  
  (2) $\Rightarrow$ (3) :
  设$\mathscr{U} = \bigcup_{i = 0}^{\infty} \mathscr{U}_i$,
  其中$\mathscr{U}_i$是局部有限(非空)开集族.
  对$\forall i \in \mathbb{N}$, $\forall U \in \mathscr{U}_i$, 记
  \[
    C(i, U) = U \setminus \bigcup_{k = 0}^{i - 1} \mathscr{U}_k,
  \]
  \[
    \mathscr{V}_i = \{ C(i, U) \mid U \in \mathscr{U}_i \}.
  \]
  其中$\mathscr{V}_0 = \mathscr{U}_0$.
  下面证$\mathscr{V} = \bigcup_{i = 0}^{\infty} \mathscr{V}_i$是$\mathscr{U}$的局部有限加细.
  对$\forall x \in X$,
  记$n(x) = \text{min}\{ i \in \mathbb{N} \mid x \in \bigcup \mathscr{U}_i \}$.
  即存在$U_x \in \mathscr{U}_{n(x)}$, 使得$x \in U_x$.
  所以有$x \in C(n(x), U_x)$, $\bigcup \mathscr{V} = X$, $\mathscr{V}$是$\mathscr{U}$的加细.
  接着, 对$\forall i > n(x)$, 有$U_x \cap \bigcup V_i = \varnothing$.
  对$\forall i \leq n(x), \exists x$的领域$A_i$,
  最多只与$\mathscr{U}_i$中的有限个元素相交,
  因此也最多只与$\mathscr{V}_i$中的有限个元素相交.
  所以$x$存在邻域$U_x \cap \bigcap_{i = 0}^{n(x)} A_i$最多和$\mathscr{V}$中的有限个元素相交.

  (3) $\Rightarrow$ (4) :
  设$\mathscr{U}$是X的一个开覆盖.
  对$\forall x \in X, \exists U_x \in \mathscr{U},
  $使得$x \in U_x, $则$x \notin U_x^c, U_x^c$为闭集.
  再由$X$满足$T_3$公理, $\exists $开集$V_x$, 使得$\overline{V_x} \subseteq U_x$.
  $\mathscr{V} = \{ V_x \mid x \in X \}$是$X$的一个开覆盖,
  而$\mathscr{V}^{'} = \{ \overline{V} \mid V \in \mathscr{V} \}$是$\mathscr{U}$的加细.
  由条件, $\mathscr{V}$有局部有限加细$\mathscr{W} = \{ W_j \mid j \in J \}$.
  因为$\mathscr{V}^{'}$是闭覆盖,
  所以$\mathscr{W}^{'} = \{ \overline{W} \mid W \in \mathscr{W} \}$
  是$\mathscr{V}^{'}$的局部有限闭加细, 从而也是$\mathscr{U}$的局部有限闭加细.
  
  (4) $\Rightarrow$ (1) :
  设$\mathscr{A}$为$X$的一个开覆盖, 由条件,
  $\mathscr{A}$存在局部有限闭加细$\mathscr{B}$.
  对$\forall x \in X, \exists x$的邻域$C_x$, 只与$\mathscr{B}$中有限个元素相交.
  记$\mathscr{C} = \{ C_x \mid x \in X \}$, 则$\mathscr{C}$为$X$的一个开覆盖.
  再由条件, $\mathscr{C}$存在局部有限闭加细$\mathscr{D}$.
  对$\forall B \in \mathscr{B},$ 
  记$B^{'} = X \setminus \bigcup \{ D \in \mathscr{D} \mid D \cap B = \varnothing \}$.
  能看出, $B^{'}$为开集且$B \subseteq B^{'}$.
  对$\forall B \in \mathscr{B}, D \in \mathscr{D}$,
  我们有$B^{'} \cap D = \varnothing$当且仅当$B \cap D = \varnothing$.
  对$\forall B \in \mathscr{B}, \exists A_B \in \mathscr{A}$,
  使得$B \subseteq A_B$. 记$\mathscr{U} = \{ B^{'} \cap A_B \mid B \in \mathscr{B} \}$.
  因为对$\forall B \in \mathscr{B}, \exists U \in \mathscr{U},$ 
  使得$B \subseteq U$, 所以$\mathscr{U}$为$X$的一个开覆盖, 
  从而$\mathscr{U}$是$\mathscr{A}$的一个开加细.
  而$\forall D \in \mathscr{D}, D$只与$\mathscr{U}$中有限个元素相交.
  加上$\mathscr{D}$是$X$的覆盖, 得出$\mathscr{U}$是局部有限的.
  所以$\mathscr{U}$是$\mathscr{A}$的局部有限开加细, 所以$X$是仿紧的.
\end{proof}

这个定理说明了, 局部有限开加细其实是一个比较强的条件, 是可以减弱的,
比如是$\sigma$-局部有限的开加细就可以了. 对于一些特定的情况下,
这个定理对于判断某些空间是不是仿紧的, 就比较方便了.

上面的条件还可以进一步减弱一些. 因为局部有限或者$\sigma$-局部有限,
同闭包保持比起来, 可能看上去还是要强一些的, 因为前者涉及到局部的性质,
而后者并没有在明处涉及到.

下面我们来看一看闭包保持(开/闭)加细和仿紧之间的联系.

\begin{thm} \label{thm:closure-preserving}
  设X是正则的Hausdorff空间,则下面各条件等价:
  \begin{enumerate}
  \item X是仿紧的;
  \item X的任意开覆盖有闭包保持开加细;
  \item X的任意开覆盖有闭包保持加细;
  \item X的任意开覆盖有闭包保持闭加细.
  \end{enumerate}
\end{thm}

该定理证明的主要部分是(4) $\Rightarrow$ (1), 而(1) $\Rightarrow$ (2),
(2) $\Rightarrow$ (3), (3) $\Rightarrow$ (4)相对来说容易一些.
为了证明(4) $\Rightarrow$ (1), 我们需要先证明几个引理.
为了方便, 我们把频繁出现的条件列出来, 避免重复说明.

\begin{condition} \label{condition:closure-preserving}
  空间$X$是正则Hausdorff空间, 且$X$的任意开覆盖都有闭包保持闭加细.
\end{condition}

\begin{lemma} \label{lemma:indexed closure-preserving}
  在条件\ref{condition:closure-preserving} 下,
  如果$\{ U_{\alpha} \}_{\alpha \in A}$是X的一个开覆盖,
  则X存在一个闭包保持闭覆盖$\{ C_{\alpha} \}_{\alpha \in A}$,
  使得$\forall \alpha \in A, C_{\alpha} \subseteq U_{\alpha}$.
\end{lemma}
\begin{proof}
  %% PROOF HERE
  由条件, $\{ U_{\alpha} \}_{\alpha \in A}$存在一个闭包保持闭加细$\mathscr{V}$.
  对$\forall V \in \mathscr{V},$ 选取$\alpha(V) \in A$, 使得$V \subseteq U_{\alpha(V)}$.
  对$\forall \alpha \in A$, 定义
  \[
    C_{\alpha} = \bigcup \{ V \in \mathscr{V} \mid \alpha(V) = \alpha \}.
  \]
  则$C_{\alpha} \subseteq U_{\alpha}$. 由$\mathscr{V}$是闭包保持闭加细,
  $C_{\alpha}$是闭集, 且$\{ C_{\alpha} \}_{\alpha \in A}$是闭包保持的.
\end{proof}

\begin{lemma} \label{lemma:regular}
  在条件\ref{condition:closure-preserving} 下, X是正规的.
\end{lemma}
\begin{proof}
  %% PROOF HERE
  设$E_1, E_2$是$X$中的两个互不相交的闭集,
  则$\{ X \setminus E_1, X \setminus E_2 \}$是X的一个开覆盖.
  由引理\ref{lemma:indexed closure-preserving},
  $X$存在闭覆盖$\{ C_1, C_2 \}$,
  使得$C_1 \subseteq X \setminus E_1, C_2 \subseteq X \setminus E_2$.
  则$E_1 \subseteq X \setminus C_1, E_2 \subseteq X \setminus C_2$,
  而$(X \setminus C_1) \cap (X \setminus C_2) = \varnothing$,
  且都是开集. 所以X是正规的.
\end{proof}

\begin{lemma} \label{lemma:Dowker}
  \text{(Dowker).} 设X是正规的, $\{ V_\gamma \}_{\gamma \in \Gamma}$是X的离散的开子集族.
  如果对$\forall \gamma \in \Gamma, \text{有}D_{\gamma} \subseteq V_{\gamma}$,
  且$\bigcup_{\gamma \in \Gamma} D_{\gamma}$是闭集,
  则X存在一个离散的开子集族$\{ G_{\gamma} \}_{\gamma \in \Gamma}$, 使得
  \[
    D_{\gamma} \subseteq G_{\gamma} \subseteq V_{\gamma}, \,\,\, \forall \gamma \in \Gamma.
  \]
\end{lemma}
\begin{proof}
  %% PROOF HERE
  记
  \[
    A = \{ x \in X \mid x \text{存在邻域至多与一个} V_{\gamma} \text{相交非空} \}.
  \]
  容易得到$A$为开集, 且
  $\bigcup_{\gamma \in \Gamma} D_\gamma \subseteq A$.
  因为$\bigcup_{\gamma \in \Gamma} D_\gamma$是闭集, $X$是正规的,
  所以存在$X$中开集$B$, 使得
  \[
    \bigcup_{\gamma \in \Gamma} D_\gamma \subseteq B \subseteq \overline{B} \subseteq A.
  \]
  令
  \[
    G_\gamma = V_\gamma \cap B.
  \]
  则$G_\gamma$是开集, 且对$\forall \gamma \in \Gamma,
  D_\gamma \subseteq G_\gamma \subseteq V_\gamma$.
  由$\{ V_\gamma \}_{\gamma \in \Gamma}$是离散的, $\{ G_\gamma \}_{\gamma \in \Gamma}$也是离散的.
\end{proof}

有了上面的准备, 我们终于可以开始这个定理的证明了.

\begin{proof}[证明(定理\ref{thm:closure-preserving})]
  %% PROOF HERE
  (4) $\Rightarrow$ (1): 设$\{ U_{\alpha} \}_{\alpha \in A}$是$X$的一个开覆盖, 且$A$是一个良序指标集.
  我们需要证明的是$\{ U_{\alpha} \}_{\alpha}$存在$\sigma$-局部有限开加细.
  
  \underline{Step 1}.
  我们先构造一列子集族, 对$\forall i \in \mathbb{N}$,
  $\{ C_{\alpha, i} \}$满足下面的条件: \newline
  \textcircled{a}
    $\{ C_{\alpha, i} \}_{\alpha \in A}$是$X$的闭包保持闭覆盖,
    且对$\forall \alpha \in A, C_{\alpha, i} \subseteq U_{\alpha}$.\newline
  \textcircled{b}
    $C_{\alpha, i+1} \cap C_{\beta, i} = \varnothing,$ 对$\forall \alpha > \beta$.\newline
  构造方法如下. 由引理\ref{lemma:indexed closure-preserving},
  存在$\{ C_{\alpha, 0} \}_{\alpha \in A}$满足\textcircled{a}(自然地满足\textcircled{b}),
  此时$i = 0$. 假设对$i \leq n, \{ C_{\alpha, i} \}_{\alpha \in A}$均满足
  \textcircled{a}和\textcircled{b}. 下面我们构造$\{ C_{\alpha, n+1} \}_{\alpha \in A}$.
  我们定义
  \[
    U_{\alpha, n+1} = U_\alpha - ( \bigcup_{\beta < \alpha} C_{\beta, n} ),\,\,\,\forall \alpha \in A.
  \]
  因为$\{ C_{\alpha, n} \}_{\alpha \in A}$是闭包保持的, 所以$U_{\alpha, n+1}$为开集.
  对$\forall x \in X$, 存在最小的$\alpha_0$, 使得$x \in U_{\alpha_0}$,
  同时, $x \in U_{\alpha_0, n+1}$, 所以$\{ U_{\alpha, n+1} \}_{\alpha \in A}$是X的一个开覆盖.
  再由引理\ref{lemma:indexed closure-preserving},
  存在一个闭包保持闭覆盖$\{ C_{\alpha, n+1} \}_{\alpha \in A}$,
  使得对$\forall \alpha \in A, C_{\alpha, n+1} \subseteq U_{\alpha, n+1}$,
  从而$C_{\alpha, n+1} \subseteq U_{\alpha}$, 因此\textcircled{a}满足.
  由$U_{\alpha, n+1}$的定义,
  有$U_{\alpha, n+1} \cap C_{\beta, n} = \varnothing, \forall \beta < \alpha$.
  因为$C_{\alpha, n+1} \subseteq U_{\alpha, n+1}$, 
  所以$C_{\alpha, n+1} \cap C_{\beta, n} = \varnothing, \forall \beta < \alpha$.
  \textcircled{b}也满足了.构造成功.

  \underline{Step 2}.
  我们继续构造一列子集族. 对$\forall \alpha \in A, i \in \mathbb{N}$, 定义
  \[
    V_{\alpha, i} = X \setminus (\bigcup_{\beta \neq \alpha} C_{\beta, i}).
  \]
  我们有\newline
  \textcircled{c}
    $\{ V_{\alpha, i} \mid \alpha \in A, i \in \mathbb{N} \}$是$X$的一个开覆盖,
    且对$\forall \alpha \in A, i \in \mathbb{N}, V_{\alpha, i} \subseteq U_\alpha$. \newline
  \textcircled{d}
    $V_{\alpha, i} \cap V_{\beta, i} = \varnothing$, 当$\alpha \neq \beta$时. \newline
  因为$\{ C_{\alpha, i} \}_{\alpha \in A}$是闭包保持的闭覆盖, 所以$V_{\alpha, i}$是开集,
  且我们有$V_{\alpha, i} \subseteq C_{\alpha, i} \subseteq U_{\alpha, i}, \forall \alpha, i$.
  再由$V_{\alpha, i}$的定义, 可知\textcircled{d}满足.
  对于\textcircled{c}, 我们还需要证明
  $\{ V_{\alpha, i} \mid \alpha \in A, i \in \mathbb{N} \}$是$X$的一个覆盖.
  任取$x \in X$, 由集合$A$是良序的, 令
  \[
    \alpha_i = \min \{ \alpha \in A \mid x \in C_{\alpha, i} \}, \,\,\, i \in \mathbb{N}.
  \]
  再选取一个自然数$k$, 使得
  \[
    \alpha_k = \min \{ \alpha_i \mid i \in \mathbb{N} \}.
  \]
  由定义, 对$\forall \alpha < \alpha_k, x \notin C_{\alpha, i}$.
  对$\forall \alpha > \alpha_k$, 由\textcircled{b},
  $C_{\alpha, k+1} \cap C_{\alpha_k, k} = \varnothing$,
  所以$x \notin C_{\alpha, k+1}$. 但由于$\{ C_{\alpha, k+1} \}_{\alpha \in A}$是$X$的一个覆盖,
  所以$x$属于$C_{\alpha_k, k+1}$. 再由$V_{\alpha_k, k+1}$的定义,
  有$x \in V_{\alpha_k, k+1}$.
  所以$\{ V_{\alpha, i} \mid \alpha \in A, i \in \mathbb{N} \}$是$X$的覆盖.

  \underline{Step 3}.
  由引理\ref{lemma:indexed closure-preserving},
  存在$X$的闭包保持闭覆盖$\{ D_{\alpha, i} \mid \alpha \in A, i \in \mathbb{N}$,
  使得对$\forall \alpha, i, D_{\alpha, i} \subseteq V_{\alpha, i}$.
  由引理\ref{lemma:regular}, X是正规的, 对$\forall i$, 由引理\ref{lemma:Dowker},
  存在离散开子集族$\{ G_{\alpha, i} \}_{\alpha \in A}$, 使得
  \[
    D_{\alpha, i} \subseteq G_{\alpha, i} \subseteq V_{\alpha, i}, \,\,\, \forall \alpha.
  \]
  于是, $\{ G_{\alpha, i} \mid \alpha \in A, i \in \mathbb{N} \}$是
  $\{ U_\alpha \}_{\alpha \in A}$的$\sigma$-局部有限开加细. 所以$X$是仿紧的.

  (1) $\Rightarrow$ (2) :
  因为$X$是仿紧的, 所以任意开覆盖存在局部有限开加细.
  由引理\ref{lemma:local finite close union}, 局部有限子集族是闭包保持的,
  所以任意开覆盖也存在闭包保持的开加细.

  (2) $\Rightarrow$ (3) :
  这是显然的.

  (3) $\Rightarrow$ (4) :
  设$\mathscr{U}$是$X$的一个开覆盖.
  因为$X$是正则的, 所以存在$X$的开覆盖$\mathscr{V}$, 使得
  $\{ \overline{V} \mid V \in \mathscr{V} \}$是$\mathscr{U}$的加细.
  由假设, $\mathscr{V}$存在闭包保持加细$\mathscr{W}$.
  而$\{ \overline{W} \mid W \in \mathscr{W} \}$是
  $\{ \overline{V} \mid V \in \mathscr{V} \}$的闭包保持闭加细,
  所以也是$\mathscr{U}$的闭包保持闭加细.
\end{proof}

到这里, 我们就已经有了好几个(在一定分离公理下)仿紧性的等价条件了.
以后判断一个空间是不是仿紧, 也就没有必要只盯着定义来回思考了,
还可以考虑这么多个等价的条件. 路多了, 总有一条适合你.


% -*- coding: utf-8 -*-
\chapter{仿紧空间的一些性质(下)}
\label{chap04}

同样是仿紧空间的性质, 和上一个性质的集合不一样,
这里的性质需要借助上面对仿紧性质的刻画才能够得到.
这也是为什么会把性质放在两个地方.

\begin{thm} \label{thm:T3 Lindelof}
  满足$T_3$公理的Lindel\"{o}f空间是仿紧的.
\end{thm}
\begin{proof}
  设$X$是满足$T_3$公理的Lindel\"{o}f空间.
  由定理\ref{thm:paracompact sigma finite},
  要证$X$仿紧,只要证$X$的任意开覆盖有$\sigma$-局部有限开加细.
  而$X$是Lindel\"{o}f的,所以任意开覆盖有可数子覆盖,
  自然是$\sigma$-局部有限开加细,所以是仿紧的.
\end{proof}

\begin{corollary}
  设拓扑空间X是$T_3$的和$C_2$的,则X是仿紧的.
\end{corollary}
\begin{proof}
  %% PROOF HERE
  因为X是$C_2$的,设$\mathscr{B}$为X的一个可数拓扑基.则对X的任意开覆盖,
  $\mathscr{B}$都是这个开覆盖的一个可数子覆盖.由此可知, $C_2$空间都是Lindel\"{o}f空间.
  再由定理\ref{thm:T3 Lindelof}, 拓扑空间X是仿紧的.
\end{proof}

上面的推论, 说明了仿紧性, 是比较普遍的, 满足了$T_3$和$C_2$即可.
而其实很多拓扑空间, 都会满足. 也就是说仿紧性, 及其可应用的范围是很广的.

我们知道, 紧空间的$F_\sigma$子空间不一定是紧空间, 但是对于仿紧性来说, 就不是这样了.
下面的定理就是说的这样一件事.

\begin{thm}
  设X是$T_3$的仿紧空间,则X的$F_{\sigma}$子空间是仿紧的.
\end{thm}
\begin{proof}
  %% PROOF HERE
  (回忆: $F_{\sigma}$子集是可数个闭子集的并.)
  设$Y$为$X$的$F_\sigma$子空间, 则有
  \[
    Y = \bigcup_{i = 1}^{\infty} A_i\,,\,\,\, A_i\text{为X中的闭集}.
  \]
  设$\mathscr{C}$为$Y$的一个开覆盖. 因为$Y$是$X$的子空间, 
  所以存在$X$的开子集族$\mathscr{C}^{'}$, 使得
  $\mathscr{C} = \{ Y \cap C \mid C \in \mathscr{C}^{'} \}$.
  接下来, 我们通过$X$的仿紧性, 来构造$\mathscr{C}$的$\sigma$-局部有限开加细.
  令
  \[
    \mathscr{W}_i = \mathscr{C}^{'} \cup \{X \setminus A_i\},
    \,\,\, i = 1, 2, \dots\,.
  \]
  易知, $\mathscr{W}_i$为$X$的开覆盖, 由$X$的仿紧性,
  $\mathscr{W}_i$存在局部有限开加细$\mathscr{F}_i$. 令
  \[
    \mathscr{G}_i = \{ F \in \mathscr{F} \mid F \cap A_i \neq \varnothing \}.
  \]
  则$\mathscr{G}_i$是$\mathscr{C}^{'}$的部分加细, 且是$A_i$的开覆盖.
  令
  \[
    \mathscr{H}_i = \{ G \cap Y \mid G \in \mathscr{G}_i \}, 
  \]
  则$\mathscr{H}_i$是$\mathscr{C}$的局部有限部分开加细. 令
  \[
    \mathscr{U} = \bigcup_{i = 1}^{\infty} \mathscr{H}_i ,
  \]
  因为$\mathscr{H}_i$是$A_i$的开覆盖, 所以$\mathscr{U}$是$Y$的开覆盖.
  因为$\mathscr{H}_i$是$\mathscr{C}$的部分开加细,
  所以$\mathscr{U}$是$\mathscr{C}$的开加细.
  而$\mathscr{U}$是$\sigma$-局部有限的, 可以得到$Y$是仿紧的.
\end{proof}

最后, 我们说明一下乘积空间仿紧性的问题.

\begin{lemma} \label{lemma:projection}
  设拓扑空间X是紧致的.对任意拓扑空间Y,投射
  \[
    p : X \times Y \longrightarrow Y,\,\, p(x, y) = y
  \]
  是完备映射.
\end{lemma}
\begin{proof}
  %% PROOF HERE
  首先, $p$是连续满映射. 对$\forall y \in Y$,
  $f^{-1}(y) = X \times \{y\}$. 因为$X$是紧的, $\{y\}$也是紧的,
  所以$f^{-1}(y)$是紧的. 接下来, 只要说明$p$是闭映射即可.
  假设$A$是$X \times Y$上的闭集, 任取$y_0 \in p(A)^{c}$,
  则有$X \times \{ y_0 \} \cap A = \varnothing$, 
  即对$\forall x \in X, (x, y_0) \notin A$.
  由$A$是闭集, 分别存在$X$和$Y$中的开集$U_x, V_x$,
  使得$(x, y_0) \in U_x \times V_x, (U_x \times V_x) \cap A = \varnothing$.
  $\{ V_x \mid x \in X \}$是$Y$的开覆盖, 由$X$的紧性,
  存在有限子覆盖$\{ U_{x_1}, U_{x_2}, \dots, U_{x_n} \}$.
  令$V = \bigcap_{i = 1}^{n} V_{x_i}$, 则$V$是$y_0$的开邻域,
  且$(X \times V) \cap A = \varnothing$,
  所以$V \cap p(A) = \varnothing$, 所以$p(A)$是闭集.
  从而$p$是完备映射.
\end{proof}

\begin{lemma} \label{lemma:perfectly paracompact}
  X, Y为拓扑空间. 设$f : X \longrightarrow Y$是完备映射, 且Y是仿紧的,
  则X是仿紧的.
\end{lemma}
\begin{proof}
  %% PROOF HERE
  设$\mathscr{U}$为$X$的开覆盖. 对$\forall y \in Y, f^{-1}(y)$是$X$的紧子集,
  从而$\mathscr{U}$存在有限子覆盖$\mathscr{U}_y$,
  使得$f^{-1}(y) \subseteq \bigcup \mathscr{U}_y$.
  由$f$的连续性, 存在$y$的邻域$V_y$,
  使得$f^{-1}(V_y) \subseteq \bigcap \mathscr{U}_y$.
  $\{ V_y \mid y \in Y \}$是$Y$的开覆盖, 由$Y$的仿紧性,
  存在局部有限开加细$\mathscr{V}$.
  对$\forall V \in \mathscr{V}, \exists y(V) \in Y$,
  使得$V \subseteq V_{y(V)}$. 我们定义
  \[
    \mathscr{A} = \{ f^{-1}(V) \cap W \mid V \in \mathscr{V}, W \in \mathscr{U}_{y(V)} \}.
  \]
  易知, $\mathscr{A}$是$\mathscr{U}$的开加细.
  因为$\mathscr{V}$是局部有限的, $f$是完备映射,
  所以$\{ f^{-1}(V) \mid V \in \mathscr{V} \}$在$X$中是局部有限的.
  而$\mathscr{U}_{y(V)}$都是有限子集族,
  所以可得$\mathscr{A}$也是局部有限的. 所以$X$是仿紧的.
\end{proof}

\begin{proposition}
  设X是紧空间, Y是仿紧空间, 则$X \times Y$是仿紧空间.
\end{proposition}
\begin{proof}
  %% PROOF HERE
  记投射
  \[
    p : X \times Y \longrightarrow Y,\,\, p(x, y) = y.
  \]
  由引理\ref{lemma:projection}, $p$是完备映射.
  因为Y是仿紧的, 由引理\ref{lemma:perfectly paracompact},
  $X \times Y$是仿紧的.
\end{proof}

%\include{body/chap05}
%\include{body/chap06}

%结论
%\include{body/chap07}

%附录
%\begin{appendix}
%   \input{appendix/appendix1.tex}
%   \input{appendix/appendix2.tex}
%  \input{appendix/appendix3.tex}
%\end{appendix}

%参考文献
% -*- coding: utf-8 -*-
\defaultfont
\chapter*{参考文献}
\markboth{参考文献}{参考文献} \addcontentsline{toc}{chapter}{参考文献}

%\defaultfont
%\fontsize{11pt}{11pt}
%\chapter{参考文献}
%\label{reference}

[1] 蒋继光. 一般拓扑学专题选讲, 四川教育出版社, 1991年

[2] 尤承业. 基础拓扑学讲义, 北京大学出版社, 1997年

[3] 儿玉之宏, 永见启应. 拓扑空间论, 2001年

[4] Ernest Michael. A Note on Paracompact Spaces(1953)

[5] Ernest Michael. Another Note on Paracompact Spaces(1956)

[6] Marco Manetti. Topology(2015) p133-135

[7] Lisa Jeffrey. Notes on Paracompactness(2015)

[8] Wikipedia. Paracompact Space

[9] Wikipedia. Cover(topology)



%\wuhao
%\clearpage
%\phantomsection
%\addcontentsline{toc}{chapter}{\bibname}

%\bibliographystyle{unsrt}
%\bibliography{body/reference}

%发表的文章列表
%\include{appendix/publications}

%致谢
% -*- coding: utf-8 -*-
\defaultfont
\chapter*{致~~~~谢}
\markboth{致谢}{致谢} \addcontentsline{toc}{chapter}{致谢}

感谢南大这四年的栽培. 感谢师维学老师在拓扑学上的教学与指导.
感谢这么多年室友对我的包容. 最后感谢家人对我的支持,
没有他们也就没有今天的我.


\clearpage
\end{document}

%%%%%%%%%%%%%%%%%% End of the file  %%%%%%%%%%%%%%%%%%%%%%%%
